\section{Context 的流动：Elys 的底层机制}

Tristan 多次强调，Elys 的核心是 context 的流动。Context 不是传统互联网里的兴趣标签，而是一个人的经历、表达方式、价值观、审美、关系、历史作品和当前意图。Elys 希望用户把这些 context 交给分身，分身再主动去社交环境中表达、评论、推荐、破冰，最终把真人连接带回来。

\begin{figure}[H]
\centering
\includegraphics[width=0.94\textwidth]{context-flow.png}
\caption{Elys 的 Context 流动闭环。}
\end{figure}

\begin{knowledgebox}{读图：Context 不是静态资料页}
图中 context 在用户、赛博分身、AI 行动、真人连接和反馈之间循环。它不是“我喜欢音乐”这种低维标签，而是能让分身像你、替你说话、替你识别关系机会的高维状态。
\end{knowledgebox}

\subsection{为什么传统互联网是低维标签}

传统互联网社交通常依赖昵称、头像、地域、年龄、兴趣、关注关系和行为日志。推荐系统能根据这些标签做分发，但很难理解一个人的价值观、幽默感、表达方式、人生经验和关系目标。长音频、长文本、聊天记录、作品和真实社交反馈在 AI 时代才开始有机会被统一理解。

\begin{importantbox}{Context 是 AI 社交的资产}
智能可能逐渐平权，但 context 不平权。谁拥有更长期、更真实、更可授权使用的用户 context，谁就更接近 AI 社交网络的核心资产。
\end{importantbox}

\subsection{Context 飞轮}

用户给出初始 context，分身开始行动；行动带来互动，互动产生新反馈；反馈又让分身更像用户。这个飞轮最难的是冷启动：用户一开始不信任产品，不愿交出隐私；产品必须在用户流失前让他看到 context 带来的真实收益。

\begin{warningbox}{Context 飞轮的脆弱点}
如果用户交出的 context 太少，分身不像自己，产品无感；如果产品索取过多，用户担心隐私和误代表。Elys 的难点正是在收益和风险之间找到启动点。
\end{warningbox}

\subsection{本章小结}

Elys 的底层不是聊天，而是 context 的循环。AI 社交是否成立，取决于 context 能否被安全获取、有效流动，并转化为真人连接。

